%% DECK-LOCAL copy of the TRM's sch_tb_pixel.tex (the system testbench),
%% reduced to the channel itself: the two 0 V current probes IPRB1 / IPRB_F
%% are plain wires, the injected accuracy-truth source I_far is removed (in
%% DPV and EIS the current comes from the cell), and the bench net names are
%% dropped.  Everything else, including the figure/resizebox wrapper the
%% \trmfit macro strips, is the TRM's.  Regenerate from the TRM fragment if
%% the topology changes.
\providecommand{\MaestroRoot}{include/analog/}
\input{\MaestroRoot sch_cell_symbol.tex}
\begin{figure}[htbp]
  \centering
  \resizebox{\linewidth}{!}{%
\ctikzset{bipoles/length=1.0cm, resistors/scale=0.85, capacitors/scale=0.85, sources/scale=0.85}
\begin{circuitikz}[
    every node/.append style={font=\small},
    nlab/.style={font=\footnotesize, inner sep=1.5pt},
    opt/.style={font=\footnotesize, text=black!55, align=center, inner sep=2pt},
    railnode/.style={circle, fill=black, inner sep=1.05pt},
    prb/.style={draw, circle, minimum size=9pt, inner sep=0pt},
  ]
  %% ---------------------------------------------------------------- amplifiers
  %% circuitikz puts `-' on top and `+' on the bottom of an `op amp'; every branch
  %% that meets a pin is placed at that pin's height with (x,0 |- pin), so no run
  %% is ever diagonal.
  \node[op amp] (A) at (0,0) {};
  \node at (A.center) {\small A1};
  \node[op amp] (B) at (13.4,0) {};
  \node at (B.center) {\small A2};

  %% ------------------------------------------------- three-electrode cell
  %% Drawn first so that (cs-ce)/(cs-re)/(cs-we) exist for the wiring below.
  \CellSymThree{6.0,0}{1.2}
  \node[opt, anchor=north] at (6.0,-1.62) {three-electrode cell};

  %% ------------------------------------------- potential program -> A1 non-inv
  \coordinate (VIP) at (-2.9,0 |- A.+);
  \draw (A.+) -- (VIP);
  \draw (VIP) to[V, l_=$v_{\mathrm{step}}$] (-2.9,-1.75);
  \draw (-2.9,-1.75) to[V, l_=$v_{\mathrm{ref}}$] (-2.9,-3.4) node[ground]{};

  %% ------------------------------ reference-electrode feedback through IPRB1
  \coordinate (VIN) at (-1.6,0 |- A.-);
  \draw (A.-) -- (VIN);
  \draw (VIN) -- (-1.6,2.6);
  \draw (-1.6,2.6) -- (6.0,2.6) -- (cs-re);

  %% ------------------------------- A1 class-AB compensation, C1/C2 back to CE
  \draw (A.up) -- (A.up |- 0,1.5);
  \draw (A.up |- 0,1.5) to[C, l=$C_c$] (2.0,1.5) -- (2.0,0);
  \node[nlab, anchor=east, xshift=-2pt] at (A.up |- 0,1.08) {\texttt{C1}};
  \draw (A.down) -- (A.down |- 0,-1.5);
  \draw (A.down |- 0,-1.5) to[C, l_=$C_c$] (2.0,-1.5) -- (2.0,0);
  \node[nlab, anchor=east, xshift=-2pt] at (A.down |- 0,-1.08) {\texttt{C2}};

  %% ------------------------------------- A1 output -> CE, WE -> the WE node
  \draw (A.out) -- (2.0,0);
  \node[railnode] at (2.0,0){};
  \draw (2.0,0) -- (cs-ce);
  \draw (cs-we) -- (10.0,0);

  %% --------------------------------------------- WE node, I_far, A2 inverting
  \node[railnode] at (10.0,0){};
  \node[nlab, anchor=north west, xshift=2pt, yshift=-2pt] at (10.0,0) {\texttt{WE}};
  \draw (10.0,0) -- (10.0,2.4);
  \draw (B.-) -- (10.0,0 |- B.-);
  \node[railnode] at (10.0,0 |- B.-){};

  %% ------------------------------------------------------ V_CM at A2 non-inv
  \draw (B.+) -- (11.7,0 |- B.+);
  \draw (11.7,0 |- B.+) to[V, l_=$V_{\mathrm{CM}}$] (11.7,-3.4) node[ground]{};

  %% ---------------------------- A2 class-AB compensation, C1/C2 back to V_OUT
  \draw (B.up) -- (B.up |- 0,1.5);
  \draw (B.up |- 0,1.5) to[C, l=$C_c$] (15.0,1.5);
  \node[nlab, anchor=east, xshift=-2pt] at (B.up |- 0,1.08) {\texttt{C1}};
  \draw (B.down) -- (B.down |- 0,-1.5);
  \draw (B.down |- 0,-1.5) to[C, l_=$C_c$] (15.0,-1.5) -- (15.0,0);
  \node[nlab, anchor=east, xshift=-2pt] at (B.down |- 0,-1.08) {\texttt{C2}};

  %% ----------------------------------- transimpedance feedback through IPRB_F
  \draw (10.0,2.4) -- (12.2,2.4);
  \node[railnode] at (12.2,2.4){};
  \draw (12.2,2.4) to[R, l=$R_f$] (15.0,2.4);
  \draw (15.0,2.4) -- (15.0,0);
  \node[railnode] at (15.0,1.5){};

  %% ------------------------------------------------------------ output node
  \draw (B.out) -- (15.0,0);
  \node[railnode] at (15.0,0){};
  \node[nlab, anchor=north west, xshift=2pt, yshift=-2pt] at (15.0,0) {$v_{\mathrm{out}}$};
  \draw (15.0,0) -- (16.2,0);
  \draw (16.2,0) to[C, l=$C_{\mathrm{out}}$] (16.2,-3.4) node[ground]{};
\end{circuitikz}
  }
  \end{figure}
